\BourbakiExercises{Exercises}

\begin{enumerate}
\def\labelenumi{\arabic{enumi})}
\item
  Prove that a left Artinian and right Noetherian ring A is also right Artinian (view the right A-modules \(\Re(A)^{k}/\Re(A)^{k+1}\) as \((A/\Re(A))\) -modules).
\item
  Let A be a (left) Artinian ring and r its radical.
\end{enumerate}

\begin{enumerate}
\def\labelenumi{\alph{enumi})}
\item
  Prove that every right ideal b of A contains a minimal right ideal (consider the ideals \(b \cap r^{p}\) , and observe that if br = 0, then b is a right A/r-module).
\item
  Prove that the left (resp. right) socle of A is the right (resp. left) annihilator of t. Deduce that every minimal two-sided ideal of A is contained in the intersection of the left and right socles of A.
\end{enumerate}

\begin{enumerate}
\def\labelenumi{\arabic{enumi})}
\setcounter{enumi}{2}
\tightlist
\item
  Let A be an Artinian ring such that the ring A/ \(\Re(A)\) is simple. Prove that A is isomorphic to a matrix ring \(\mathbf{M}_{r}(\mathrm{B})\) over an Artinian local ring B (use Exercises 19 and 20 of VIII, p.~169, as well as Exercise 19 of VIII, p.~94).
\end{enumerate}

¶ 4) A ring A is called pseudoregular if for every element a of A, there exist an integer n and an element x of A such that \(a^{n}xa^{n}=a^{n}\) . An absolutely flat ring (VIII, p.~171, Exercise 27) is pseudoregular.

\begin{enumerate}
\def\labelenumi{\alph{enumi})}
\item
  Prove that a pseudoregular ring is a Zorn ring (VIII, p.~171, Exercise 26) in which every noninvertible element is a left and right zero divisor.
\item
  Prove that every quotient ring of a pseudoregular ring is pseudoregular (cf.~VIII, p.~171, Exercise 27, d)).
\item
  Prove that the center Z of a pseudoregular ring A is pseudoregular (prove that if \(a \in Z\) and \(x \in A\) satisfy \(a^{n}xa^{n} = a^{n}\) , then \(a^{2n}x^{k}\) belongs to Z for every \(k \geqslant 1\) ).
\item
  Let p be a prime number, \(U_{p}\) the p-primary component of the group Q/Z, and A the (commutative) algebra of the group \(U_{p}\) over the field \(F_{p}\) . Prove that the ring A is pseudoregular but \(A^{N}\) is not.
\end{enumerate}

¶ 5) Let n be an integer. A ring A is called n-regular if for every element a of A, there exists an \(x \in A\) such that \(xa^{n+1} = a^n\) .

\begin{enumerate}
\def\labelenumi{\alph{enumi})}
\item
  Suppose that \(\mathrm{A}\) is \(n\) -regular. Prove that we have \(r^n = 0\) for every element \(r\) of \(\Re(\mathrm{A})\) .
\item
  Prove that an n-regular primitive ring is isomorphic to a matrix ring \(\mathbf{M}_{r}(\mathrm{D})\) over a field, with \(r \leqslant n\) . (Observe that if V is a D-vector space of dimension \textgreater{} n and e a nonzero vector in V, then every dense subring of \(\operatorname{End}_{\mathrm{D}}(\mathrm{V})\) contains an element u such that \(u^{n}(e) \neq 0\) and \(u^{n+1}(e) = 0\) .)
\item
  Prove that an n-regular ring is a Zorn ring (VIII, p.~171, Exercise 26). (First, observe that the relation \(xa^{n+1}=a^{n}\) implies \(x^{n}a^{2n}=a^{n}\) . Using b), Exercise 12, a) of VIII, p.~168 and Proposition 2 of VIII, p.~27, deduce that we have \(a^{n}x^{n}a^{n}-a^{n}\in\Re(A)\) . Finally, using Lemma 1 of VIII, p.~159, prove that if two elements y and z of A satisfy \(y=zy^{2}\) and \(yzy-y\in\Re(A)\) , then there exists a \(t\in A\) such that \(y=ty^{2}\) and \((ty)^{2}=ty.\) If the ring A is, moreover, without radical, then it is pseudoregular (Exercise 4).
\end{enumerate}

\(d\) ) Let A be a ring with no nilpotent elements other than 0. Then A is \(n\) -regular if and only if it is absolutely flat.

(If A is absolutely flat, use Exercise 27, b) of VIII, p.~171. If A is n-regular, first observe, with the notation above, that \(e = x^{n}a^{n}\) is idempotent, hence central (Exercise 30, a) of VIII, p.~172), and that we have ea = a, so that A is 1-regular. Then use questions a) and b) to prove that A is isomorphic to a subring of a product of fields and deduce that the relation \(xa^{n+1} = a^{n}\) implies axa = a.)

¶ 6) Let A be a Artinian ring and \(n = \text{long}_{\text{A}}(\text{A}_s)\) .

\begin{enumerate}
\def\labelenumi{\alph{enumi})}
\item
  Prove that the ring A is \(n\) -regular (Exercise 5).
\item
  Prove that the ring A is pseudoregular. (Let r be an integer such that \(\Re(A)^{r}=0\) , and let m be the length of the ring B = A/ \(\Re(A)\) . Prove, by reducing to the case when the ring B is simple, that for every \(b\in B\) , there exists an element y of B that commutes with b and satisfies \(b^{m+1}y=b^{m}\) . Deduce that for every \(a\in A\) , there exists an \(x\in A\) such that \(a^{mr}xa^{mr}=a^{mr}\) .)
\item
  Let \(\mathfrak{I}\) be the set of nonnilpotent left ideals of A. Prove that the minimal elements of \(\mathfrak{I}\) are the ideals \(Ae\) generated by an indecomposable idempotent (use \(a\) ) and Exercise 5, \(c\) ). Prove that every left ideal of A is the direct sum of finitely many ideals of this type and a left ideal contained in the radical.
\item
  A left ideal a of A is generated by an idempotent if and only if it is the direct sum of a finite family of left ideals \((Ae_{i})_{i\in I}\) , where the \(e_{i}\) are indecomposable orthogonal idempotents.
\end{enumerate}

¶ 7) Let A be an Artinian ring, r its radical, and \(\overline{A}\) the ring A/r. Let l be a nonnilpotent indecomposable left ideal; we have l = Ae, where e is an idempotent in A.

\begin{enumerate}
\def\labelenumi{\alph{enumi})}
\item
  Prove that the ideal \(\mathfrak{n} = \mathfrak{l} \cap \mathfrak{r}\) is the unique maximal submodule of \(\mathfrak{l}\) . The A-module \(\mathfrak{l} / \mathfrak{n}\) is simple and isomorphic to \(\overline{\mathrm{A}}\overline{e}\) , where \(\overline{e}\) is the image of \(e\) in \(\overline{\mathrm{A}}\) .
\item
  Let M be an A-module. Then \((Ae)M\) is a submodule of M, and \(e\big((Ae)M\big)=eM\) is an eAe-module. If N is an A-submodule of M, then the eAe-modules \(e(M/N)\) and eM/eN are isomorphic. If the A-module M is simple, then either we have eM=0, or M is isomorphic to \(\overline{A}\overline{e}\) and eM is a simple eAe-module.
\item
  Prove that the length of the eAe-module eM is equal to number of quotients of a Jordan--Hölder series of the A-module M isomorphic to \(\overline{A}\overline{e}\) .
\item
  Suppose that M has finite length and admits a unique maximal submodule N and that M/N is isomorphic to \(l/n\) . Prove that M is isomorphic to a quotient of Ae (observe that an element x of eM that does not belong to N generates M).
\end{enumerate}

¶ 8) Let A be an Artinian ring. The A-module \(A_{s}\) is the direct sum of a finite family of indecomposable left ideals \((Ae_{i})_{i\in I}\) , where the \(e_{i}\) are orthogonal idempotents (Exercise 6, d)).

\begin{enumerate}
\def\labelenumi{\alph{enumi})}
\tightlist
\item
  Two A-modules of finite length M and N are said to be linked if a quotient of a Jordan--Hölder series of M is isomorphic to a quotient of a Jordan--Hölder series of N. Let R be the finest equivalence relation on the set of ideals \(Ae_{i}\) for which two linked ideals are equivalent. The equivalence classes \(\mathcal{B}_{l}\) ( \(l \in L\) ) for R are called blocks; we denote the sum of the ideals in block \(B_{l}\) by \(b_{l}\) . Prove that the ideals \(b_{l}\) are the only indecomposable two-sided ideals of A.
\end{enumerate}

(To see that \(b_{l}\) is two-sided, observe that we have \(Ae_{i}xe_{j} \subset Ae_{j}\) and that by Exercise 7, c), the relation \(e_{i}Ae_{j} \neq 0\) implies \(Ae_{i} \equiv Ae_{j} \pmod{R}\) . On the other hand, let \((\mathfrak{b}_{m}^{\prime})_{m \in \mathrm{M}}\) be a finite family of two-sided ideals with direct sum A. Prove that each of the \(b_{m}^{\prime}\) is a sum of ideals \(Ae_{i}\) and then that two linked ideals \(Ae_{r}\) and \(Ae_{s}\) belong to the same ideal \(b_{m}^{\prime}\) by observing that by Exercise 7, there exists an index h such that \(e_{h}Ae_{r}\) and \(e_{h}Ae_{s}\) are \(\neq 0\) .)

\begin{enumerate}
\def\labelenumi{\alph{enumi})}
\setcounter{enumi}{1}
\item
  Prove that the right ideals \(e_{i}A\) are indecomposable (consider their images in A/ \(\Re(A)\) ).
\item
  Suppose that A is an algebra of finite degree over a commutative field K. We denote by \(c_{ij}\) (resp. \(d_{ij}\) ) the number of quotients of a Jordan--Hölder series of \(Ae_{j}\) (resp. \(e_{j}A\) ) isomorphic to \(\overline{A}\overline{e}_{i}\) (resp. \(\overline{e}_{i}\overline{A}\) ) and by \(r_{i}\) the degree of the field \(e_{i}Ae_{i}/e_{i}\Re(A)e_{i}\) over K. Prove the equality \(c_{ij}r_{j}=d_{ij}r_{i}\) (determine the dimension of \(e_{i}Ae_{j}\) over K, and use Exercise 7, c)).
\end{enumerate}

\begin{enumerate}
\def\labelenumi{\arabic{enumi})}
\setcounter{enumi}{8}
\tightlist
\item
  Let A be an Artinian ring whose radical \(\mathfrak{r}\) is contained in the center.
\end{enumerate}

\begin{enumerate}
\def\labelenumi{\alph{enumi})}
\item
  Prove that every central idempotent in A/r lifts to a central idempotent in A (if e is an idempotent in A with central class in A/r, express the fact that e commutes with xe - ex for every \(x \in A\) ).
\item
  If \(\mathfrak{r}\) is nonzero, prove that \(A / \mathfrak{r}\) is commutative (observe that we have the relation \((xy - yx)z = 0\) for \(x\in A\) , \(y\in A\) , and \(z\in \mathfrak{r}\) ).
\item
  Deduce from a) and b) that A is isomorphic to the product of a semisimple ring and a finite family of local rings \(A_{i}\) such that the field \(\mathrm{A}_{i}/\Re(\mathrm{A}_{i})\) is commutative. d) Prove that every simple A-module is isomorphic to a minimal left ideal of A.
\end{enumerate}

¶ 10) For a ring A, denote the set of pseudosubrings of A whose elements are all nilpotent by \(\mathcal{N}(\mathrm{A})\) and the set of maximal elements of \(\mathcal{N}(\mathrm{A})\) by \(\mathcal{N}_{\max}(\mathrm{A})\) .

\begin{enumerate}
\def\labelenumi{\alph{enumi})}
\tightlist
\item
  Let D be a field and n an integer. Prove that \(\mathcal{N}_{\max}(\mathbf{M}_{n}(\mathrm{D}))\) consists of the pseudorings \(gTg^{-1}\) , where g is an invertible element of \(\mathbf{M}_{n}(\mathrm{D})\) and T is the set of matrices \((a_{ij})\in\mathbf{M}_{n}(\mathrm{D})\) such that \(a_{ij}=0\) for \(i\geqslant j\) (use Exercise 2 of VIII, p.~14).
\item
  Let A be an Artinian ring and \(\pi:A\to A/\Re(A)\) the canonical mapping. Show that the mapping \(P\mapsto\pi(P)\) defines a bijection from \(\mathcal{N}_{\max}(A)\) to \(\mathcal{N}_{\max}(A/\Re(A))\) . Deduce from a) that the intersection of the elements of \(\mathcal{N}_{\max}(A)\) is equal to \(\Re(A)\)
\end{enumerate}

and that two arbitrary elements of \(\mathcal{N}_{\mathrm{max}}(\mathrm{A})\) are transformed into each other by an inner automorphism of A.

\begin{enumerate}
\def\labelenumi{\arabic{enumi})}
\setcounter{enumi}{10}
\tightlist
\item
  Let A be a commutative ring, B an A-algebra, and M an Artinian B-module. Suppose that the module M is finitely generated.
\end{enumerate}

\begin{enumerate}
\def\labelenumi{\alph{enumi})}
\item
  Prove that the ring \(B_{M}\) is left Artinian (use Proposition 6 of VIII, p.~9) and then that the B-module M has finite length.
\item
  Prove that if M is a simple B-module, then the A-module M is isotypical (deduce from Lemma 1 of VIII, p.~81 that the kernel of the homomorphism \(A \rightarrow \text{End}_{B}(M)\) is a maximal ideal).
\item
  Conclude that M is an A-module of finite length.
\end{enumerate}

